\documentclass[11pt]{book}
\usepackage[paperwidth=145mm,paperheight=200mm,top=18mm,bottom=20mm,inner=20mm,outer=16mm]{geometry}
\usepackage{brumfiel}

% ---- local macros (hoist into book preamble) ----
% Halftone photographs in the original are reproduced as PLATES: a framed
% placeholder carrying the original caption. Never redrawn.
%   #1 = placeholder height, #2 = caption block (already formatted)
\newcommand{\FMplate}[2]{%
  \par\medskip
  \begin{center}
  \fbox{\begin{minipage}[c][#1][c]{0.68\linewidth}\centering\footnotesize\itshape
    halftone photograph\\ not reproduced in this edition
  \end{minipage}}%
  \par\smallskip
  {\footnotesize #2}
  \end{center}
  \par\medskip}

% Title-page author line: name in letterspaced caps, affiliation in italic.
\newcommand{\FMauthor}[2]{%
  \par\medskip
  {\large\scshape #1}\par
  {\footnotesize\itshape #2}\par}
% Narrow measure (145 mm trim) + Palatino: give TeX a little slack so the
% justified prose paragraphs set without overfull lines.
\emergencystretch=1.5em

% Blank versos forced by \chapter* (openright) must carry no folio or head.
\makeatletter
\def\cleardoublepage{\clearpage
  \if@twoside\ifodd\c@page\else
    \null\thispagestyle{empty}\newpage
    \if@twocolumn\hbox{}\newpage\fi
  \fi\fi}
\makeatother
% ---- end local macros ----

\begin{document}
\pagenumbering{roman}

%========================================================= half title
\thispagestyle{empty}
\begin{center}
  \vspace*{0.28\textheight}
  {\Huge GEOMETRY}
\end{center}
\clearpage

%========================================================= p. i  series page
\thispagestyle{empty}
\null\vfill
\begin{center}
  {\footnotesize This book is in the}\par\smallskip
  {\normalsize ADDISON-WESLEY SCIENCE EDUCATION SERIES}\par
  \rule{0.28\linewidth}{0.4pt}\par\smallskip
  {\footnotesize\itshape Consulting Editors}\par\smallskip
  {\footnotesize
   \begin{tabular}{@{}ll@{}}
   \textsc{Richard S. Pieters} \quad & \textsc{Paul Rosenbloom}\\
   \textsc{George B. Thomas, Jr.} \quad & \textsc{John Wagner}
   \end{tabular}}
\end{center}
\vfill
\clearpage

%========================================================= p. ii  frontispiece
% Euclid frontispiece: engraved/halftone portrait plate facing the title page.
\thispagestyle{empty}
\null\vspace*{0.06\textheight}
\FMplate{52mm}{%
  \textsc{Euclid}\\
  Courtesy of\\
  \textit{Scripta Mathematica}}
\vfill
\clearpage

%========================================================= p. iii  title page
\thispagestyle{empty}
\null\vspace*{0.10\textheight}
\begin{center}
  {\Huge GEOMETRY}\par
  \vspace{2.2em}
  \FMauthor{Charles F. Brumfiel}{Ball State Teachers College}
  \FMauthor{Robert E. Eicholz}{Ball State Teachers College, Burris Laboratory School}
  \FMauthor{Merrill E. Shanks}{Purdue University}
  \vspace{3.2em}

  {\normalsize ADDISON-WESLEY PUBLISHING COMPANY, INC.}\par
  {\footnotesize READING, MASSACHUSETTS, U.S.A.}\par
  {\footnotesize LONDON, ENGLAND}
\end{center}
\vfill
\clearpage

%========================================================= p. iv  copyright
% Original copyright statement, transcribed as printed in the 1961 second
% printing. Reproduced as a historical record of the source edition.
\thispagestyle{empty}
\null\vspace*{0.22\textheight}
\begin{center}
  {\itshape Copyright \textcopyright\ 1960}\par
  {\normalsize ADDISON-WESLEY PUBLISHING COMPANY, INC.}\par
  \rule{0.30\linewidth}{0.4pt}\par
  {\itshape Printed in the United States of America}\par
  \vspace{1.1em}
  {\footnotesize\textsc{All rights reserved. This book, or parts thereof,
   may not be reproduced in any form without written permission of the
   publisher.}}\par
  \vspace{0.9em}
  {\itshape Library of Congress Catalog Card No. 60-8336}\par
  \vspace{1.4em}
  {\normalsize Second printing, February, 1961}
\end{center}
\vfill
\clearpage

%========================================================= contents
% The printed book's CONTENTS list (book pp. v--viii) is deliberately NOT
% transcribed: the assembled edition generates its own table of contents from
% the \addcontentsline entries in each unit. The integration agent places
% \tableofcontents here.

%========================================================= PREFACE  (p. ix)
\chapter*{\centering\Large PREFACE}
\addcontentsline{toc}{chapter}{Preface}
\markboth{PREFACE}{PREFACE}

It is a common belief that plane geometry was completed by Euclid 2000 years ago
and that nothing has been added to it or taken from it since. This is simply not
true. That there are logical gaps in Euclid's presentation has been known for a
long time. Means to remedy these deficiencies have been known for about sixty
years, but strangely enough a mathematically adequate and yet elementary
treatment of plane geometry in the spirit of Euclid has not appeared in print.
This text represents an earnest effort to do just this. The current interest in
improvement of the secondary school curriculum makes this an appropriate time
for such an attempt.

This book, in earlier unpublished versions, has been taught to \emph{normal}
classes in \emph{average} high schools for five years. Many of the classes were
taught by one of us in the Burris School, the laboratory school of Ball State
Teachers College, but the material has been taught in other schools by teachers
who had had no previous special training in geometry. Our experience is that
both student and teacher find this treatment of geometry stimulating. The good
students become excited about mathematics and the students of average ability
learn the basic geometric facts.

Besides its mathematical completeness, an attribute that makes for clarity and
ease of understanding, there are several features which we think important.

There is a full chapter on logic early in the text. It pays the teacher to
return to this logic again and again in order to clarify the principles
underlying deductive proof. For the nonscience student, this chapter on logic
may be the item of greatest permanent value.

The postulates are presented in groups, as needed, with a resulting gain in
clarity and simplicity. Although it is true that at the start most students are
hardly ``provers,'' by the end of the year a considerable part of the class will
exhibit real mathematical perception. In the early stages the small number of
postulates makes proofs easier to construct.

Definitions are given explicitly; that is, they are formally labeled as such. It
is one of our prime objectives to get the student to see that precise
definitions are necessary for proofs.

There is plenty of material for a full year's work for all kinds of classes.
More difficult sections, theorems, or exercises are marked with an asterisk, and
some of the more difficult theorems may be treated as postulates. The chapters
on analytic geometry, locus, solid geometry, and philosophy of mathematics
provide supplementary material beyond the standard course. The chapter on solid
geometry alone could occupy a good portion of a semester if desired.

There is a review of important concepts at the end of each chapter. In addition,
there are algebra review sections, designed not only to refresh skills but also
to bring out the logical structure of algebra.

The Appendix lists all the postulates, plus the principal definitions and
theorems. Besides being a convenient reference, the Appendix is worth studying
during the year, for it gives, in brief, the step-by-step construction of the
full geometry.

We believe that the text raises students' reading ability. There is little
verbiage. Each section has a definite purpose, and the text is directed to this
purpose. Good classroom discussion, with careful attention to the use of words,
will add much to each student's ability to express himself clearly.

A separate Teachers' Manual gives greater detail on selection of topics,
emphasis, and time schedule. There is also an Answer Book for the exercises.

It is a great pleasure to thank all those who have helped us with encouragement
or constructive criticism. Special thanks are due to Dr.\ Curtis Howd, Principal
of the Burris Laboratory School, who encouraged us and made possible the
experimental sections; to Professor Jack Forbes, who observed many of the
classes in various schools and who insisted that the mathematics be correct; to
teachers at the Burris School, and many others, who by teaching the material
have improved the exposition; to The National Science Foundation, which, through
a grant for in-service institutes at Ball State, enabled us to reach a wider
audience and see the text taught by many teachers in a variety of school
situations; and finally, to our wives, who have waited impatiently for the
publication day.

\vspace{1.2em}
\noindent
\begin{minipage}[t]{0.45\linewidth}
\itshape January 1960
\end{minipage}\hfill
\begin{minipage}[t]{0.45\linewidth}
\raggedleft C.\,B., R.\,E., M.\,S.
\end{minipage}

%========================================================= FOREWORD  (p. xi)
\chapter*{\centering\Large FOREWORD TO THE STUDENT}
\addcontentsline{toc}{chapter}{Foreword to the Student}
\markboth{FOREWORD TO THE STUDENT}{FOREWORD TO THE STUDENT}

Mathematics is much more than a collection of interesting facts and useful tools
for solving special problems. Mathematics is primarily concerned with the
\emph{development} of important ideas. Frequently mathematicians develop an
interesting mathematical theory that is apparently useless so far as the
solution of practical problems is concerned. Later someone finds that this
mathematics fits very nicely into the pattern of some complicated practical
problem. The mathematical theory then becomes a useful tool.

As a branch of mathematics, geometry is first of all concerned with the
development of geometrical theory. As you study geometry you will see again and
again how the geometrical concepts fit into patterns in the real world and
become helpful in the solution of practical problems.

Mathematics requires thought, and thinking deserves your best effort. Do not
hesitate to ask questions.

\vspace{1.0em}
\noindent
\begin{flushright}
Good luck,\par
\textit{The Authors}
\end{flushright}

%========================================================= p. xii  Hilbert plate
% Halftone photograph on book p. xii, facing the Chapter 1 opener (book p. 1).
% Kept here because p. xii falls inside the front-matter page range; if the
% ch01 unit also carries it, drop one copy at integration.
\clearpage
\thispagestyle{empty}
\null\vspace*{0.05\textheight}
\FMplate{58mm}{%
  \textsc{David Hilbert}\\
  \textit{(1862--1943)}\\
  Professor of Mathematics,\\
  University of G\"ottingen, Germany\\[2pt]
  Perhaps the world's foremost mathematician during the\\
  first half of the twentieth century.\\[2pt]
  \textit{Courtesy of the New York Public Library.}}
\vfill

\end{document}
